\begin{exemplo}[Fully connected registers: internal degeneracy allowed]
\label{ex:cliques}
For $k\ge3$, take $k$ modules $V_0,\dots,V_{k-1}$ of $m=4$ vertices each, with
$n=4k$. Within each module, coupling $1$ on \emph{all} six edges ---
the module induces $K_4$. A bridge edge of weight $w=0.10$ joins the last vertex of
$V_i$ to the first vertex of $V_{i+1}$. The site energies equal $\Delta_{\mathrm{mod}}\,i$
at all four vertices of $V_i$, with $\Delta_{\mathrm{mod}}=24$.

The spectrum of the block is $\{24i+3,24i-1\}$, with multiplicities $1$ and $3$
\cite{brouwer2012}. Thus the family preserves internal degeneracies
at all sizes.

\paragraph{Verification of (H1)--(H3)}
The bridge edges are pairwise disjoint, since bridge $i$ is incident to the last
vertex of $V_i$ and to the first vertex of $V_{i+1}$, which are distinct. Hence $\Hout$ is a direct
sum of $2\times2$ blocks of the form
$\left(\begin{smallmatrix}0&w\\w&0\end{smallmatrix}\right)$ and
$\normop{\Hout}=w=0.10=:\delta$, which gives (H1). The clusters are
$\Sigma_i=\{24i-1,\ 24i+3\}$, of width $W_i=4$, and consecutive clusters are at distance
$(24(i+1)-1)-(24i+3)=20$; hence
\[
g=\Delta_{\mathrm{mod}}-m=20\ >\ 2\delta=0.20,
\]
which gives (H2). Each module induces a clique, so $D_i=1$ and (H3) holds with $D_0=1$.
The maximum degree is $d=m=4$ (three internal neighbors, plus at most one bridge) and
$\mu=\max\{1,w\}=1$. Consequently
\[
\eta=\frac{0.10}{19.90}\approx0.0050251,
\qquad
\ghat=\frac{20-0.20}{8}=2.4750,
\]
\[
\varrho_i=1+\frac{2(4+0.20)}{\pi\cdot19.80}\approx1.135041 .
\]

\paragraph{Bounds}
Choose $s$ to be the first vertex of $V_0$ and $t$ the last vertex of
$V_{k-1}$. The structure of cliques joined by a single bridge gives, exactly,
\[
\dist(s,V_i)=2i,\qquad
\dist(t,V_i)=2(k-1-i),\qquad
D=\dist(s,t)=2k-1.
\]
\begin{table}[t]
\centering
\small
\setlength{\tabcolsep}{4pt}
\caption{Certificates for the fully connected registers of
Example~\ref{ex:cliques} ($K_4$ modules, bridges of weight $0.10$, detuning
$24$), from the first vertex of $V_0$ to the last vertex of $V_{k-1}$. The
columns are described at the beginning of Section~\ref{sec:exemplos}.}
\label{Tab:cliques}
\begin{tabular}{rrrrlll}
\toprule
$k$ & $n$ & $D$ & $\kappa_{\rm tab}$ & $\Sigma_1\le$ & $F_{\rm esp}\le$ & $B^2\le$\\
\midrule
$3$ & $12$ & $5$ & $0.5000$ & $0.01008$ & $1.016\times10^{-4}$ & $1.011\times10^{-4}$\\
$4$ & $16$ & $7$ & $0.5000$ & $0.01011$ & $1.021\times10^{-4}$ & $1.011\times10^{-4}$\\
$5$ & $20$ & $9$ & $0.8440$ & $3.588\times10^{-3}$ & $1.288\times10^{-5}$ & $1.288\times10^{-5}$\\
$6$ & $24$ & $11$ & $0.8712$ & $4.279\times10^{-4}$ & $1.831\times10^{-7}$ & $1.831\times10^{-7}$\\
$8$ & $32$ & $15$ & $0.8615$ & $2.353\times10^{-5}$ & $5.535\times10^{-10}$ & $5.535\times10^{-10}$\\
$10$ & $40$ & $19$ & $0.8681$ & $8.805\times10^{-7}$ & $7.752\times10^{-13}$ & $7.752\times10^{-13}$\\
\bottomrule
\end{tabular}
\end{table}

\noindent The subspace bound, independent of $k$, already yields
$\Fmax\le b_\eta^2\le1.011\times10^{-4}$ throughout the family.
At $k=5$ ($n=20$), the spatial certificate $\Sigma_1$ improves this value to
$\Fmax\le1.288\times10^{-5}$, and the improvement becomes more pronounced in the
following rows of Table~\ref{Tab:cliques}.
The factor $\varrho_i\approx1.135$ measures the cost of enclosing the full interval
and remains close to $1$ in this family.

\end{exemplo}

\begin{exemplo}[Detuned dimerized chain]\label{ex:dimero}
For $k\ge3$, take the path with $n=2k$ vertices grouped into $k$ modules
$V_i=\{2i,2i+1\}$: internal coupling $J=1$ on the edge of each module, coupling
$w=0.05$ on the $k-1$ bridge edges, and site energies $H_{2i,2i}=24i$ and
$H_{2i+1,2i+1}=24i+10$.

The bridge edges form a matching, so $\Hout$ is a direct sum of
blocks $\left(\begin{smallmatrix}0&w\\w&0\end{smallmatrix}\right)$ and
$\normop{\Hout}=w=0.05=: \delta$. Each block
$H_i=\left(\begin{smallmatrix}24i&1\\1&24i+10\end{smallmatrix}\right)$ has
$\Sigma_i=\{24i+5\pm\sqrt{26}\}$, of width
$W_i=2\sqrt{26}\approx10.198039$, and consecutive clusters are at distance
$g=24-2\sqrt{26}\approx13.801961>2\delta=0.10$. With $d=2$, $\mu=1$,
$D_0=1$ and $\abs{V_i}=2$:
\[
\eta\approx0.0036358,\qquad \ghat\approx3.425490,\qquad
\varrho_i\approx1.478467 .
\]
For $s=0$ and $t=2k-1$,
\[
\dist(s,V_i)=2i,\qquad
\dist(t,V_i)=2(k-1-i),\qquad
D=2k-1.
\]

\begin{table}[t]
\centering
\small
\setlength{\tabcolsep}{4pt}
\caption{Certificates for the detuned dimerized chain of
Example~\ref{ex:dimero} (bridges of weight $0.05$), from $s=0$ to $t=2k-1$.
Columns as in Table~\ref{Tab:cliques}.}
\label{Tab:dimero}
\begin{tabular}{rrrrlll}
\toprule
$k$ & $n$ & $D$ & $\kappa_{\rm tab}$ & $\Sigma_1\le$ & $F_{\rm esp}\le$ & $B^2\le$\\
\midrule
$3$ & $6$ & $5$ & $0.5000$ & $7.285\times10^{-3}$ & $5.307\times10^{-5}$ & $5.288\times10^{-5}$\\
$5$ & $10$ & $9$ & $0.8565$ & $5.477\times10^{-4}$ & $3.000\times10^{-7}$ & $3.000\times10^{-7}$\\
$7$ & $14$ & $13$ & $0.8830$ & $1.721\times10^{-5}$ & $2.960\times10^{-10}$ & $2.960\times10^{-10}$\\
$9$ & $18$ & $17$ & $0.8739$ & $1.983\times10^{-7}$ & $3.931\times10^{-14}$ & $3.931\times10^{-14}$\\
$12$ & $24$ & $23$ & $0.9006$ & $9.191\times10^{-11}$ & $8.448\times10^{-21}$ & $8.448\times10^{-21}$\\
\bottomrule
\end{tabular}
\end{table}

\noindent The subspace bound is $b_\eta^2\le5.288\times10^{-5}$,
independently of $k$. At $n=10$, the spatial certificate $\Sigma_1$ of
Table~\ref{Tab:dimero} already yields $\Fmax\le3.000\times10^{-7}$. The rate is governed by the interval separation
$g\approx13.802$. The internal spacing $2\sqrt{26}$ influences the width
and the value of $g$, but does not appear as an additional simplicity hypothesis.

\end{exemplo}

\begin{exemplo}[$P_4$ modules: larger communities]\label{ex:p4}
For $k\ge3$, take the path with $n=4k$ vertices grouped into $k$ modules of four
consecutive vertices, each inducing a $P_4$: coupling $1$ on the three internal
edges, $w=0.3$ on the $k-1$ bridges, and site energy $5i$ at all vertices
of $V_i$. Again the bridges form a matching, so
$\normop{\Hout}=w=0.3=: \delta$. With
$\phi=(1+\sqrt5)/2$, the known spectrum of the path~\cite{brouwer2012} gives
\[
\spec(P_4)=\{\pm\phi,\ \pm\phi^{-1}\},\qquad
W_i=2\phi=1+\sqrt5,
\]
and therefore
\[
g=5-W_i=4-\sqrt5\approx1.763932>2\delta=0.6.
\]
With $d=2$, $\mu=1$, $D_0=3$ and $\abs{V_i}=4$:
\[
\eta\approx0.2049275,\qquad \ghat\approx0.290983,
\qquad \varrho_i\approx3.098161 .
\]
For the endpoint vertices $s=0$ and $t=4k-1$,
\[
\dist(s,V_i)=4i,\qquad
\dist(t,V_i)=4(k-1-i),\qquad
D=4k-1.
\]

\begin{table}[t]
\centering
\small
\setlength{\tabcolsep}{4pt}
\caption{Certificates for the $P_4$ modules of Example~\ref{ex:p4} (bridges of
weight $0.3$), from $s=0$ to $t=4k-1$. Columns as in Table~\ref{Tab:cliques}.}
\label{Tab:p4}
\begin{tabular}{rrrrlll}
\toprule
$k$ & $n$ & $D$ & $\kappa_{\rm tab}$ & $\Sigma_1\le$ & $F_{\rm esp}\le$ & $B^2\le$\\
\midrule
$10$ & $40$ & $39$ & $0.8642$ & $0.2474$ & $0.06118$ & $0.06118$\\
$20$ & $80$ & $79$ & $0.8951$ & $7.449\times10^{-4}$ & $5.548\times10^{-7}$ & $5.548\times10^{-7}$\\
$40$ & $160$ & $159$ & $0.9453$ & $1.290\times10^{-11}$ & $1.662\times10^{-22}$ & $1.662\times10^{-22}$\\
\bottomrule
\end{tabular}
\end{table}

\noindent Here $\varrho_i\approx3.10$, since the internal width is
approximately $3.24$ and the residual gap approximately $1.16$.
The subspace bound is $b_\eta^2\le0.1610$.
Because $\Sigma_1$ keeps the minima with $b_i$ that the expanded
form~\eqref{eq:principalexpandida} drops, the spatial certificate is informative
already in the row $n=40$ of Table~\ref{Tab:p4}, with $\Fmax\le0.06118$. At
$n=80$, it gives $\Fmax\le5.548\times10^{-7}$.
\end{exemplo}
